\section{The three boundary means}\label{app:boundary} In \(\Gamma\) each square has been isolated, but
its two diagonal pairs can be available simultaneously. To preserve the sign and count their interaction
once, we start with the correction for one pair and then combine both pairs in the square. Here \(p_e\)
denotes the stationary \emph{rook} occupancy on the binary mask. For an available diagonal pair
\(\parmar=\{e,f\}\), let \(\pi^{\rm ex}_{\parmar}\) be its joint probability in a uniform rook placement.
Set \(\pi^T_{\parmar}=\exp[L_T(M_{\parmar})-L_T(M)]\), where \(M_{\parmar}\) deletes its four row and
column endpoints. With \(\mathsf G(e,f)\) the cross response of the stationary rook point and \(\chi_e\)
its single-square relative correction, define \[
 \pi_{1,\parmar}=p_ep_f(1-\mathsf G(e,f)+\chi_e+\chi_f),\qquad
 X_{\parmar}=p_ep_f(p_e+p_f)\mathsf G(e,f).
\] For a single pair the square correction is \(\log(1-\pi^{\rm
ex}_{\parmar})-\log(1-\pi_{1,\parmar})-X_{\parmar}+O(j^{-5})\). The probabilities are \(O(j^{-2})\) and
their difference is \(O(j^{-4})\). Integrating the logarithmic derivative gives a linearization error of
\(O(j^{-6})\). The full square contributes the corrections for its two pairs and an additional term
\(-3\prod_{s\in\square}p_s+O(j^{-5})\). There are at most \(j^2(j-1)\) pairs, since each square has at
most \(2(j-1)\) neighbors and each pair has two endpoints. Thus the following quantity \(\mathcal
M_n(b)\) differs from \(n\mathbb
 E_\mu[\Gamma;F_b]\) by \(O(n/j^2)=o_b(1)\). \begin{align*}
 \mathcal M_n(b)={}&-n\mathbb E_\mu[\textstyle\sum_{\parmar}(\pi^{\rm ex}_{\parmar}-\pi^T_{\parmar});F_b]\\
           &-n\mathbb E_\mu[\textstyle\sum_{\parmar}(\pi^T_{\parmar}-\pi_{1,\parmar}+X_{\parmar});F_b]
            -3n\mathbb E_\mu[\textstyle\sum_\square\prod_{s\in\square}p_s;F_b].
\end{align*} Introducing \(\pi^T_{\parmar}\) in this decomposition serves two purposes. The first mean
measures the error of the rook predictor on the pair. The second retains the insertion error for a marked
pair. To average it, we need to know which of the other two vertices remain available. Let \(u_f\) be the
final vacancies of the two diagonal families, and set \(U_f(b)=u_f+b(1-u_f)\). Let \(v_b=U_D(b)U_A(b)\),
let \(p_b\) be its doubly stochastic scaling~\cite{sk}, and put \(W_b^0=p_b/v_b\). A square with minimum
coordinates \((r,c)\) and side length \(h\) is integrated over \(\Omega=\{0<h<1,\ 0<r,c<1-h\}\). With
\(d_s=U_D(b,r-c+sh)\) and \(a_s=U_A(b,r+c+sh)\), the availability of its oriented pair contributes the
factor \(B_o\). The other two vertices have indicators \(I_l I_u,I_l I_v\). Their four joint
probabilities and the factor for the marked pair are \[
\begin{array}{c|cccc}
o&B_o&l&u&v\\\hline
 D&d_0a_0a_2&a_1&d_{-1}&d_{+1}\\
 A&a_1d_{-1}d_{+1}&d_0&a_0&a_2
\end{array}
\qquad
\begin{array}{c|c}
s&q_{o,s}\\\hline
00&1-l+l(1-u)(1-v)\\
10&lu(1-v)\\
01&l(1-u)v\\
11&luv
\end{array}
\] Each state has submeasure \(b^4 B_o q_{o,s}\,dr\,dc\,dh\). The four live labels are two rows and two
columns. A repeated diagonal vacancy retains the same indicator. Each orientation has
\(\sum_{h=1}^{n-1}(n-h)^2\) positions. The Riemann normalization contributes \(n^3\), with four factors
of \(j^{-1}\) and an outer factor of \(n\). Hence \(n^4 b^4/j^4\to1\). Inserting the two marked weights
gives \(w_{o,b}=B_oW_b^0(e)W_b^0(f)=p_b(e)p_b(f)/U_o(b)\). Even states of zero mass are retained, with no
division by \(q_{o,s}\).

For the first mean, let \(m_{f,b}\) be the mass of \(p_b\) on each line of family \(f\), with the same
fiber measure \(dr\) used for the final mass \(m_f\) in Section~\ref{sec:profile-data}. Set \[
 W_b=\sum_{f=D,A}\int\frac{m_{f,b}^2}{2U_f},\qquad
 \gamma_b=\frac{m_{D,b}}{U_D}+\frac{m_{A,b}}{U_A},\qquad
 \theta_b(r)=\int p_b(r,c)\gamma_b(r,c)\,dc .
\] \paragraph{The increment of a double minor.} Fix an admissible cut-off \(b>0\) and \(j=\lfloor
bn\rfloor\). On a regular binary mask \(M\) of size \(j\), write \(A_R,A_C\) for the uniform measures on
the live rows and columns. Let \(\mathfrak F(M)\) be the rook functional in Section~\ref{sec:torres}.
When differentiating it in its measures we retain the input mask and recompute its scaling so that both
margins remain equal to one. Write
\[
 h_R^M(i)=D_R\mathfrak F(M)[\delta_i-A_R],\qquad
 h_C^M(l)=D_C\mathfrak F(M)[\delta_l-A_C].
\]
Each influence has mean zero with respect to its own measure, by linearity. For the pair
\(\parmar=\{e,f\}\) write \(H_M(\parmar)=\sum_{z=e,f}(h_R^M(r_z)+h_C^M(c_z))\). In the minor
\(M_{\parmar}\) the measures assign mass \(1/(j-2)\) to each remaining line. In particular,
\[
 A'_R-A_R=\frac{2A_R-\delta_{r_e}-\delta_{r_f}}{j-2},\qquad
 A'_C-A_C=\frac{2A_C-\delta_{c_e}-\delta_{c_f}}{j-2}.
\]
This is an exact measure identity. Put \(\Delta_{\parmar}\mathfrak F=\mathfrak F(M_{\parmar})-\mathfrak
F(M)\). The uniform bound on the second derivative in these directions (Cauchy's estimate in the masses
of the distinguished lines) gives, for every pair in the class,
\[
 \Delta_{\parmar}\mathfrak F=-\frac{H_M(\parmar)}{j-2}+O(j^{-2}),\qquad
 j\Delta_{\parmar}\mathfrak F=-H_M(\parmar)+O(j^{-1}).
\]
Let \(D\) be the discrete doubly stochastic scaling of \(M\), and let \(w_{\parmar}=(jD_e)(jD_f)\), which
is bounded in the class. The ratio of the exact pair probability to its rook approximation and the
approximation itself satisfy
\[
 \frac{\pi_{\parmar}^{\rm ex}}{\pi_{\parmar}^T}=\exp\{g_T(M_{\parmar})-g_T(M)\},\qquad
 \pi_{\parmar}^T=j^{-2}w_{\parmar}(1+O(j^{-1})).
\]
Indeed, the first equality divides the ratios of the permanents and of their predictors. The second uses
the leading term of the two-site marginal and \(p_z/D_z=1+O(j^{-1})\). Upon subtracting
\(\coefT(M_{\parmar})/(j-2)-\coefT(M)/j\), with \(\coefT=\mathfrak F-1/12\), the two successive
increments of the remainder of \((T^+)\) are \(O(j^{-3})\). Thus,
\[
 j^2\{g_T(M_{\parmar})-g_T(M)\}
 =j\Delta_{\parmar}\mathfrak F+2\mathfrak F(M)-1/6+O(j^{-1}).
\]
Expanding the exponential in the exact ratio then gives
\begin{equation}\label{eq:pt-par-fisico}
 -j^4(\pi_{\parmar}^{\rm ex}-\pi_{\parmar}^T)
 =w_{\parmar}\{1/6-2\mathfrak F(M)-j\Delta_{\parmar}\mathfrak F\}+O(j^{-1}).
\end{equation}
The remainder is uniform over pairs. There are at most \(j^2(j-1)\) of them, and the first mean in
\(\mathcal M_n(b)\) multiplies this sum by \(n/j^4\). The total cost of this remainder and of replacing
\(-j\Delta_{\parmar}\mathfrak F\) by \(H_M(\parmar)\) is therefore
\[
 O\!\left(\frac{n}{j^4}\,j^2(j-1)\,\frac1j\right)
 =O(n/j^2)=o_b(1).
\]
 \paragraph{The density of the removed lines.} The first mean in \(\mathcal M_n(b)\) sums over the states
of both crossings and \(\sum_s q_{o,s}=1\). The pattern law and the scaling at regular arrival, developed
in Appendix~\ref{app:transfer}, give the pair measure \(w_{o,b}\,dr\,dc\,dh\) for each orientation.
Indeed, the four axial labels contribute \(b^4\), each distinct diagonal contributes a vacancy, and the
two factors \(jD_z\) converge to \(W_b^0(z)\) on the support. The outer factor \(n/j^4\) and the \(n^3\)
from the Riemann sum leave \(n^4 b^4/j^4\to1\). By symmetry between the endpoints of the pair, on each
line with coordinate \(t\) its mass integrates to \(\frac12\iint p_b(z)p_b(z')\,dr\,dr'/U_o(t)\). Summing
over lines gives \(W_b\).

If each pair contributes a test function \(a\) at each of its two rows, the symmetry of \(a(r)+a(r')\)
cancels this factor \(1/2\) when one integrates a single endpoint. In the family \(o\), integrating the
other endpoint gives \(m_{o,b}(t)/U_o(t)\). Thus,
\[
 \sum_o\int_\Omega w_{o,b}\{a(r_e)+a(r_f)\}\,dr\,dc\,dh
 =\iint a(r)p_b(r,c)\gamma_b(r,c)\,dr\,dc
 =\int a(r)\theta_b(r)\,dr .
\]
Transposition gives the same density for columns. In particular, \(\int\theta_b=2W_b\). The row measure
lost to the pairs has density \(\theta_b\), and renormalizing the surviving rows restores the mass
\(2W_b\) at constant density. Together the two operations give the direction \(\eta_b=2W_b-\theta_b\).

The centering can be verified exactly on a deterministic grid with \(m\) points per side. Let \(p_m\) be
the kernel of \(v_b\) scaled with the uniform measure \(A_m\) of mass \(1/m\) per point, and let
\(W_m,\theta_m\) be the grid sums corresponding to \(W_b,\theta_b\). On a line \(\ell\) of family \(o\)
we give each ordered pair the weight \(\omega=p_m(e)p_m(f)/(2m^3 U_o(\ell))\), including the pairs
\(e=f\) in each family. Summing over both endpoints gives \(\sum\omega=W_m\). Fixing the row \(i\) of the
first endpoint and multiplying by two gives \(\theta_m(i)/m\). With \(h_R^m,h_C^m\) denoting the centered
influences on this grid and \(H_m(\parmar)\) their sum over the four endpoints, linearity gives
\[
 \sum\omega H_m(\parmar)
 =\frac1m\sum_i\theta_m(i)(h_R^m(i)+h_C^m(i))
 =-D\mathfrak F_m[\eta_m A_m,\eta_m A_m],\qquad \eta_m=2W_m-\theta_m .
\]
The equality holds for the first variations. The pairs \(e=f\) in this quadrature have total mass
\(O_b(m^{-1})\). The physical minors always use two distinct squares.

\begin{proposicion}[First mean at fixed cut-off]\label{prop:pt-corte}
Keep \(v_b\) fixed and evaluate \(\mathfrak F\) with the two measures \((1+\varepsilon\eta_b)\,dr\),
\((1+\varepsilon\eta_b)\,dc\), rescaling \(v_b\) at each \(\varepsilon\). Denote the derivative at zero
by \(dF_2(b)\). Then
\begin{align}
 \lim_{n\to\infty}\frac n{j^4}\mathbb E_\mu
 \Big[\sum_{\parmar} w_{\parmar}\,j\Delta_{\parmar}\mathfrak F\,;F_b\Big]&=dF_2(b),
 \label{eq:pt-menores-limite}\\
 P_T(b):=\lim_{n\to\infty}-n\mathbb E_\mu
 \Big[\sum_{\parmar}(\pi_{\parmar}^{\rm ex}-\pi_{\parmar}^T)\,;F_b\Big]&=(1/6-2F(b))W_b-dF_2(b).
 \label{eq:pt-frontera-fija}
\end{align}
\end{proposicion}
\begin{proof}
To replace \(H_M(\parmar)\) by the profile influences we must also transport the derivative of the
scaling. In the series for the functional the centered kernel \(\centrado=jD-1\) has gap
\(\|\centrado\|_{2\to2}\le r_b<1\) and sections of norm at most \(C_b\). A walk of length \(k\ge2\)
between fixed endpoints is bounded by \(C_b^2r_b^{k-2}\). With \(T_R=(I-\centrado\centrado^*)^{-1}-I\),
the diagonal of \(T_R-\centrado\centrado^*\) already excludes the unique return of length two.
For columns the roles of \(\centrado\) and \(\centrado^*\) are interchanged. The truncated terms of the
functional are therefore sums over simple graphs. Differentiating a measure fixes a vertex and subtracts
its mean. Differentiating the scaling in a row \(i\) adds a path to that root, since the column response
in the centered gauge is \(-(I-\centrado^*\centrado)^{-1}\centrado_{i,\cdot}\). The choices of
differentiated vertex or edge multiply the geometric tail by a polynomial in the length, so it remains
summable.

At fixed truncation, accidental coincidences of labels cost \(O_{b,k}(1/n)\). The pattern law on live
lines, uniform in the root, takes each sum to the same graph with kernel \(p_b\). At fixed precision the
exceptional event on \(F_b\) has exponentially small mass, and the scaled weights are bounded. We take
\(n\to\infty\) and then remove the truncation. Thus in the weighted mean \(\mathfrak F(M),h_R^M,h_C^M\)
are replaced by \(F(b),h_R,h_C\), with error \(o_b(1)\). Here \(h_R,h_C\) are the continuous influences,
with mean zero. Integrating over the four endpoints gives their contribution
\[
 \int h_R\theta_b\,dr+\int h_C\theta_b\,dc
 =-D\mathfrak F[\eta_b\,dr,\eta_b\,dc]=-dF_2(b).
\]
This proves~\eqref{eq:pt-menores-limite} by \(j\Delta_{\parmar}\mathfrak F=-H_M(\parmar)+O(j^{-1})\).
Substitution into~\eqref{eq:pt-par-fisico} and use of the mass \(W_b\)
prove~\eqref{eq:pt-frontera-fija}.
\end{proof}
Identity~\eqref{eq:pt-frontera-fija} holds for the limit of physical arrival on the event \(F_b\).
If the four
axial labels are fixed and the three diagonals are required to be empty in the complete solution, the
difference between the two means is \(O(b)+o_b(1)\). Uniform positivity of \(v_b\) and continuity of the
scaling and of its first variation give \(P_T=\lim_{b\downarrow0}P_T(b)\). In this term the line measures
are varied. In \(F'(0)\), which enters the interior defect, the kernel \(v_b\) is differentiated along
the profile with both measures fixed.
 \paragraph{The third mean.} For
the third mean, the square uses six distinct diagonal labels. Two of them each occur twice among the
eight vacancy factors in \(\prod_s W_b^0(s)\). Multiplying by the six distinct vacancies and setting
\(p_0=\nu_1\) leaves exactly \[
 -S=-3\int_\Omega
       \frac{\prod_{s\in\square}\nu_1(s)}
            {u_D(r-c)\,u_A(r+c+h)}\,dr\,dc\,dh .
\] This identifies the first and third means. The second retains two pieces of mask data at the marked
square, the indicators of its two intersections. Appendix~\ref{app:weights} computes the pair error with
these two indicators fixed.
